% ============================================================
% ADD THIS AFTER "Figure \ref{fig:demo_gallary}" paragraph
% ============================================================

\medskip

{\noindent\textbf{Dataset Complexity Analysis}}: To contextualize SimBench's complexity, we provide a quantitative comparison with established code generation benchmarks in Table~\ref{tab:benchmark_comparison}. Token counts are computed using OpenAI's \texttt{tiktoken} library with the \texttt{cl100k\_base} encoding.

\begin{table}[h!]
    \centering
    \caption{Comparison of SimBench with existing code generation benchmarks. SimBench exhibits significantly higher complexity in both prompt and solution lengths.}
    \label{tab:benchmark_comparison}
    \begin{tabular}{l r r r r r}
        \toprule
        \textbf{Benchmark} & \textbf{Tasks} & \textbf{Prompt} & \textbf{Solution} & \multicolumn{2}{c}{\textbf{Ratio to SimBench}} \\
        & & \textbf{(tokens)} & \textbf{(tokens)} & Prompt & Solution \\
        \midrule
        MBPP~\cite{austin2021program} & 500 & 16 & 58 & 61.6$\times$ & 24.4$\times$ \\
        MBPP+~\cite{liu2024your} & 378 & 19 & 40 & 51.8$\times$ & 35.4$\times$ \\
        DS-1000~\cite{lai2023ds} & 1,000 & 282 & 42 & 3.5$\times$ & 33.7$\times$ \\
        HumanEval~\cite{chen2021evaluating} & 164 & 131 & 54 & 7.5$\times$ & 26.2$\times$ \\
        HumanEval+~\cite{liu2024your} & 164 & 131 & 54 & 7.5$\times$ & 26.2$\times$ \\
        BigCodeBench~\cite{zhuo2024bigcodebench} & 1,140 & 145 & 112 & 6.8$\times$ & 12.6$\times$ \\
        CodeContests~\cite{li2022competition} & 165 & 593 & 762 & 1.7$\times$ & 1.9$\times$ \\
        \midrule
        \textbf{SimBench (ours)} & \textbf{102} & \textbf{985} & \textbf{1,415} & 1.0$\times$ & 1.0$\times$ \\
        \bottomrule
    \end{tabular}
\end{table}

A key distinction of SimBench lies in its multi-turn structure. Unlike single-turn benchmarks where each task is independent, SimBench's second and third turns require the S-LLM to modify existing code from the previous turn. This design reflects real-world software development workflows where developers iteratively refine and extend existing codebases. Table~\ref{tab:turn_complexity} details the prompt composition across turns.

\begin{table}[h!]
    \centering
    \caption{Prompt complexity across SimBench turns. Turns 2 and 3 include full code context from previous turns, reflecting realistic code modification scenarios.}
    \label{tab:turn_complexity}
    \begin{tabular}{l r r r r r}
        \toprule
        \textbf{Turn} & \textbf{Tasks} & \textbf{Text} & \textbf{Code Context} & \textbf{Total Prompt} & \textbf{Solution} \\
        & & \textbf{(tokens)} & \textbf{(tokens)} & \textbf{(tokens)} & \textbf{(tokens)} \\
        \midrule
        Turn 1 & 34 & 87 & --- & 87 & 1,252 \\
        Turn 2 & 34 & 189 & 1,252 & 1,441 & 1,416 \\
        Turn 3 & 34 & 123 & 1,306 & 1,428 & 1,577 \\
        \midrule
        \textbf{Average} & \textbf{102} & \textbf{133} & \textbf{853} & \textbf{985} & \textbf{1,415} \\
        \bottomrule
    \end{tabular}
\end{table}

The multi-turn paradigm introduces several challenges absent in conventional benchmarks:
\begin{itemize}
    \item \textbf{Long-context reasoning}: With prompts averaging 985 tokens (and exceeding 1,400 tokens in later turns), S-LLMs must maintain coherence across extended contexts while understanding the existing codebase.
    \item \textbf{Code comprehension}: Unlike generation-only tasks, turns 2 and 3 require the S-LLM to first understand the provided code before making targeted modifications, testing both reading and writing capabilities.
    \item \textbf{Incremental complexity}: Solutions grow progressively longer (1,252 $\rightarrow$ 1,416 $\rightarrow$ 1,577 tokens), requiring the S-LLM to manage increasing code complexity while maintaining consistency with prior implementations.
\end{itemize}

Table~\ref{tab:category_complexity} provides a breakdown of complexity by simulation category, revealing that Sensor (SEN) and Finite Element Analysis (FEA) scenarios demand the most extensive code generation.

\begin{table}[h!]
    \centering
    \caption{Dataset statistics by simulation category.}
    \label{tab:category_complexity}
    \begin{tabular}{l r r r r}
        \toprule
        \textbf{Category} & \textbf{Systems} & \textbf{Tasks} & \textbf{Avg. Prompt} & \textbf{Avg. Solution} \\
        & & & \textbf{(tokens)} & \textbf{(tokens)} \\
        \midrule
        Multibody (MBS) & 5 & 15 & 965 & 1,395 \\
        FEA & 5 & 15 & 1,275 & 1,775 \\
        Sensor (SEN) & 4 & 12 & 1,270 & 1,781 \\
        Robotics (RBT) & 6 & 18 & 800 & 1,178 \\
        Vehicle (VEH) & 14 & 42 & 887 & 1,291 \\
        \midrule
        \textbf{Total} & \textbf{34} & \textbf{102} & \textbf{985} & \textbf{1,415} \\
        \bottomrule
    \end{tabular}
\end{table}

% ============================================================
% REFERENCES TO ADD TO YOUR BIBLIOGRAPHY
% ============================================================
% @article{austin2021program,
%   title={Program synthesis with large language models},
%   author={Austin, Jacob and others},
%   journal={arXiv preprint arXiv:2108.07732},
%   year={2021}
% }
% @article{chen2021evaluating,
%   title={Evaluating large language models trained on code},
%   author={Chen, Mark and others},
%   journal={arXiv preprint arXiv:2107.03374},
%   year={2021}
% }
% @article{liu2024your,
%   title={Is your code generated by chatgpt really correct? rigorous evaluation of large language models for code generation},
%   author={Liu, Jiawei and others},
%   journal={NeurIPS},
%   year={2024}
% }
% @article{lai2023ds,
%   title={DS-1000: A natural and reliable benchmark for data science code generation},
%   author={Lai, Yuhang and others},
%   journal={ICML},
%   year={2023}
% }
% @article{zhuo2024bigcodebench,
%   title={BigCodeBench: Benchmarking Code Generation with Diverse Function Calls and Complex Instructions},
%   author={Zhuo, Terry Yue and others},
%   journal={arXiv preprint arXiv:2406.15877},
%   year={2024}
% }
% @article{li2022competition,
%   title={Competition-level code generation with alphacode},
%   author={Li, Yujia and others},
%   journal={Science},
%   year={2022}
% }
